\documentclass[11pt,a4paper]{article}
% !TEX root = ../report_plan/report_draft/RENJIEQU_235_FinalReport.tex
\usepackage[utf8]{inputenc}
\usepackage[T1]{fontenc}
\usepackage{amsmath,amssymb,amsthm}
\usepackage{geometry}
\usepackage{hyperref}
\usepackage{enumitem}
\usepackage{titlesec}
\usepackage{booktabs}
\usepackage{xcolor}
\usepackage{tcolorbox}
\usepackage{array}
\usepackage{longtable}

\geometry{margin=2.4cm}
\hypersetup{colorlinks=true, linkcolor=blue!60!black, citecolor=blue, urlcolor=blue!70!black}

% ---- bra-ket and operator macros (physics.sty not available) ----
\newcommand{\ket}[1]{\left|#1\right\rangle}
\newcommand{\bra}[1]{\left\langle#1\right|}
\newcommand{\braket}[2]{\left\langle#1\middle|#2\right\rangle}
\newcommand{\ev}[1]{\left\langle#1\right\rangle}
\newcommand{\dd}{\mathrm{d}}
\newcommand{\Tr}{\operatorname{Tr}}
\newcommand{\Real}{\operatorname{Re}}
\newcommand{\Imag}{\operatorname{Im}}
\newcommand{\half}{\tfrac{1}{2}}

\newtheorem{definition}{Definition}[section]
\newtheorem{remark}{Remark}[section]
\newtheorem{keyidea}{Key idea}[section]

% status tags
\newcommand{\known}{\textcolor{green!55!black}{\textbf{[course-covered]}}}
\newcommand{\newx}{\textcolor{red!70!black}{\textbf{[NEW -- focus here]}}}
\newcommand{\partly}{\textcolor{orange!85!black}{\textbf{[partly new]}}}

\newtcolorbox{checkbox}[1]{colback=blue!4,colframe=blue!45!black,title=#1,fonttitle=\bfseries}
\newtcolorbox{whybox}{colback=yellow!8,colframe=orange!60!black,title=Why the workflow needs this,fonttitle=\bfseries}

\titleformat{\part}[display]{\bfseries\Large\color{blue!40!black}}{\partname~\thepart}{0.4em}{\Huge}

\title{\bfseries Physics Theory Roadmap \& Study Guide\\[4pt]
\large Reproducing Quench Dynamics of the Schwinger Model\\ via Variational Quantum Algorithms}
\author{Companion to: R.~Qu, \emph{Physics 235 Final Project} (Spring 2026)\\
Source paper: Nagano, Bapat \& Bauer, \emph{Phys. Rev. D} \textbf{108}, 034501 (2023) [arXiv:2302.10933]}
\date{\today}

\begin{document}
\maketitle

\begin{tcolorbox}[colback=red!5,colframe=red!60!black,title=\textbf{Before you start: a note on your source files}]
Your proposal is built on \textbf{Nagano--Bapat--Bauer}, ``Quench dynamics of the Schwinger model
via variational quantum algorithms'' (arXiv:2302.10933). The PDF currently in your project folder
(\texttt{s41467-025-57580-5.pdf}) is a \emph{different} paper --- Bosse, Childs \emph{et al.},
``Efficient and practical Hamiltonian simulation from time-dependent product formulas''
(Nat.\ Commun.\ 2025). That paper is genuinely useful for the \emph{Trotter} part of your project
(Part IV.B below), but it is not your project's source. \textbf{Action: download arXiv:2302.10933}
and keep both. This guide is anchored to the workflow your proposal specifies.
\end{tcolorbox}

\tableofcontents
\newpage

% ==================================================================
\part{Orientation --- the whole project on one page}
% ==================================================================

\section{What you are actually building}
Your project is a pipeline that takes a continuum quantum field theory (1+1D quantum
electrodynamics, the \emph{Schwinger model}) and turns it into something a 4-qubit quantum
computer can simulate, then runs two kinds of real-time dynamics on it and checks them against an
exact classical answer. Read this pipeline top to bottom --- every later part of the guide expands
one box.

\begin{center}
\renewcommand{\arraystretch}{1.35}
\begin{tabular}{>{\raggedright\arraybackslash}p{0.30\textwidth} >{\raggedright\arraybackslash}p{0.40\textwidth} l}
\toprule
\textbf{Stage} & \textbf{What happens} & \textbf{Status} \\
\midrule
1. Choose the physics & Schwinger model = 1+1D U(1) gauge theory of one Dirac fermion & \newx \\
2. Discretize & Kogut--Susskind staggered fermions on $N{=}4$ sites; gauge field on links & \partly \\
3. Remove the gauge field & Solve Gauss's law in 1+1D $\Rightarrow$ non-local fermionic $H$ & \newx \\
4. Encode in qubits & Jordan--Wigner $\Rightarrow$ 4-qubit Pauli sum & \known \\
5. Ground state & VQE: variational minimization of $\ev{H}$; metric $r(E)$ & \known \\
6. Quench & Suddenly switch on external field $L{=}5$; state is now out of equilibrium & \partly \\
7a. Dynamics (baseline) & 2nd-order Suzuki--Trotter evolution $e^{-iHt}$ & \known \\
7b. Dynamics (paper) & McLachlan variational time evolution: $A\dot\theta = C$ & \newx \\
8. Observe \& validate & $\ev{L(t)},\ \ev{n(t)}$, fidelity $F(t)$ vs.\ exact diagonalization & \known \\
\bottomrule
\end{tabular}
\end{center}

\begin{remark}
The two ingredients your proposal flags as ``beyond course coverage'' are exactly the two
\newx{} rows that are conceptually heaviest: (i) the \textbf{gauge structure} (stages 1 and 3) and
(ii) \textbf{McLachlan variational quantum simulation} (stage 7b). This guide spends most of its
ink there and keeps the \known{} rows short, since you already have them from class.
\end{remark}

\section{Theory dependency map}
This is the table to pin above your desk. Each row is a physics theory; the columns tell you what
it feeds and how much of it is new to you (you reported: solid on QM and quantum computing, new to
gauge theory).

\begin{center}
\renewcommand{\arraystretch}{1.3}
\begin{longtable}{>{\raggedright\arraybackslash}p{0.27\textwidth} >{\raggedright\arraybackslash}p{0.45\textwidth} l}
\toprule
\textbf{Core theory} & \textbf{Feeds which stage} & \textbf{For you} \\
\midrule
\endhead
Gauge theory \& local symmetry & Stage 1, 3 (why a gauge field exists at all) & \newx \\
Schwinger model phenomenology & Stage 1, 8 (confinement, pair production --- what to expect) & \newx \\
Lattice gauge theory / staggered fermions & Stage 2 (the lattice $H$) & \partly \\
Gauss's law \& constraint elimination & Stage 3 (removing the gauge dof in 1+1D) & \newx \\
Jordan--Wigner transformation & Stage 4 (fermions $\to$ qubits) & \known \\
Variational principle / VQE & Stage 5 (ground state) & \known \\
Trotter / Suzuki product formulas & Stage 7a (baseline dynamics) & \known \\
Time-dependent variational principle (McLachlan) & Stage 7b (the paper's method) & \newx \\
Quantum geometric tensor / Fubini--Study metric & Stage 7b ($A$ matrix) & \newx \\
Hadamard-test circuits & Stage 7b (measuring $A$, $C$ on hardware) & \partly \\
Quench / out-of-equilibrium dynamics & Stage 6 (the protocol) & \partly \\
\bottomrule
\end{longtable}
\end{center}

\begin{checkbox}{Prerequisite self-check (you should be comfortable with all of these)}
\begin{itemize}[leftmargin=1.4em,nosep]
\item Second quantization: creation/annihilation operators $\chi^\dagger_n,\chi_n$, anticommutators.
\item Pauli matrices $X,Y,Z$, the spin-$\half$ algebra, tensor products of qubits.
\item Time-independent perturbation / variational principle $\ev{H}\ge E_0$.
\item Quantum circuits, parameterized gates $e^{-i\theta P/2}$, expectation-value measurement.
\item Trotter splitting $e^{-i(A+B)t}\approx (e^{-iA t/r}e^{-iB t/r})^r$.
\end{itemize}
If any of these are shaky, patch them \emph{first} --- the new material assumes them fluently.
\end{checkbox}

% ==================================================================
\part{Physics foundations: gauge theory and the Schwinger model}
% ==================================================================

\section{Gauge theory in five minutes (the part that is genuinely new)}\label{sec:gauge}

\begin{whybox}
You cannot build the Schwinger Hamiltonian without understanding \emph{why} there is an electric
field $L_n$ living on the links and \emph{why} it is constrained by Gauss's law. Everything in
stages 1--3 of the pipeline is gauge theory. This is the single highest-leverage concept to learn.
\end{whybox}

\subsection{Global vs.\ local symmetry}
A free Dirac fermion has a \emph{global} $U(1)$ symmetry: $\psi \to e^{i\alpha}\psi$ with $\alpha$
a constant leaves the Lagrangian invariant. Promoting $\alpha \to \alpha(x)$ (a different phase at
each spacetime point) is a \emph{local} or \emph{gauge} symmetry. The ordinary derivative
$\partial_\mu \psi$ is no longer covariant because $\partial_\mu(e^{i\alpha(x)}\psi)$ produces an
extra $i(\partial_\mu\alpha)\psi$ term.

\begin{keyidea}
To restore invariance you must introduce a new field $A_\mu$ --- the gauge field --- and replace
the derivative by the \emph{covariant derivative}
\begin{equation}
D_\mu = \partial_\mu + i g A_\mu, \qquad A_\mu \to A_\mu - \tfrac{1}{g}\partial_\mu\alpha .
\label{eq:covderiv}
\end{equation}
The gauge field is not optional decoration: \emph{demanding local symmetry forces it to exist}, and
it is what we call the photon / electric field. This is the logic behind ``the $U(1)$ gauge
interaction giving rise to the full Schwinger Hamiltonian'' in your proposal.
\end{keyidea}

The gauge field gets its own kinetic term from the field strength
$F_{\mu\nu}=\partial_\mu A_\nu-\partial_\nu A_\mu$, giving the QED Lagrangian
\begin{equation}
\mathcal{L} = \bar\psi(i\gamma^\mu D_\mu - m)\psi - \tfrac{1}{4}F_{\mu\nu}F^{\mu\nu}.
\label{eq:qedlagrangian}
\end{equation}

\subsection{The Schwinger model: 1+1D QED}
The Schwinger model is Eq.~\eqref{eq:qedlagrangian} in \emph{one} space + one time dimension. Two
facts make 1+1D special and make this a perfect quantum-simulation testbed:

\begin{enumerate}[leftmargin=1.5em]
\item \textbf{The gauge field is not dynamical on its own.} In 1+1D there are no transverse
photon polarizations. The electric field $E$ is fully determined by the charges through Gauss's law
$\partial_x E = g\,j^0$. This is what lets you \emph{eliminate} the gauge field entirely in stage 3
--- a luxury unavailable in 3+1D.
\item \textbf{It is still rich.} Despite being ``simple,'' the model exhibits genuinely
field-theoretic phenomena:
\begin{itemize}[nosep]
\item \emph{Confinement}: the potential between a static charge--anticharge pair grows
\emph{linearly} with separation (you cannot isolate a single charge).
\item \emph{Schwinger pair production}: a strong electric field spontaneously creates
particle--antiparticle pairs --- the dynamical effect your $\ev{n(t)}$ plot will show after the
quench.
\item \emph{Chiral condensate} $\ev{\bar\psi\psi}\neq 0$ and a dynamically generated photon mass
$m_\gamma = g/\sqrt{\pi}$ in the massless case (the original Schwinger result).
\end{itemize}
\end{enumerate}

\begin{remark}[What to physically expect in your results]
When you switch on the external field $L=5$ at the quench, you are driving the vacuum hard. Expect
$\ev{n(t)}$ (particle density) to \emph{rise} as pairs are produced, and $\ev{L(t)}$ (electric
field) to oscillate as the produced charges partially screen the field --- a real-time picture of
the Schwinger mechanism. Knowing this in advance is your sanity check on the numerics.
\end{remark}

\subsection{The continuum Hamiltonian (temporal gauge)}
Fixing $A^0=0$ (temporal gauge), the Hamiltonian density has three pieces you will recognize on the
lattice: a fermion kinetic term, a mass term, and the electric-field energy $\tfrac12 E^2$,
\begin{equation}
H = \int \dd x\left[ -i\bar\psi\gamma^1(\partial_1 + igA_1)\psi \;+\; m\,\bar\psi\psi \;+\; \tfrac12 E^2\right],
\qquad \partial_1 E = g\,\psi^\dagger\psi .
\label{eq:contH}
\end{equation}
The constraint on the right is Gauss's law; it is the seed of stage 3. Hold onto the three-term
structure --- it survives discretization almost unchanged.

% ==================================================================
\part{From field theory to a 4-qubit Hamiltonian}
% ==================================================================

\section{Kogut--Susskind staggered fermions}\label{sec:ks}

\begin{whybox}
This is stage 2: how Eq.~\eqref{eq:contH} becomes a finite lattice model you can put on qubits. You
have seen ``staggered fermions'' and ``tight-binding chain'' in class, so this section is
\partly{} --- the new piece is the \emph{gauge field on links}.
\end{whybox}

\subsection{The fermion-doubling problem and the staggered cure}
Naively discretizing a Dirac fermion produces spurious extra particles (``doublers'') --- the
lattice momentum has unwanted low-energy modes at the Brillouin-zone edge. Kogut and Susskind's fix
is to put a \emph{single-component} fermion field $\chi_n$ on each site and distribute the Dirac
spinor's components across alternating sites: even sites carry ``particles,'' odd sites carry
``antiparticles.'' This is encoded by the \emph{staggered mass term} with its alternating sign
$(-1)^n$.

\subsection{The lattice Hamiltonian}
Gauge fields live on the \emph{links} between sites. Define on link $n$ the unitary link operator
$U_n=e^{iag A_n}$ and its conjugate \emph{electric-field operator} $L_n$ with integer eigenvalues
$\ell\in\mathbb{Z}$ (compact $U(1)$), satisfying $[\,\theta_n, L_m]=i\delta_{nm}$ where
$U_n=e^{i\theta_n}$. The Kogut--Susskind Hamiltonian is
\begin{equation}
\boxed{\;
H = -\frac{i}{2a}\sum_{n}\left[\chi^\dagger_n U_n \chi_{n+1} - \text{h.c.}\right]
\;+\; m\sum_n (-1)^n \chi^\dagger_n\chi_n
\;+\; \frac{a g^2}{2}\sum_n L_n^2 \; }
\label{eq:ksH}
\end{equation}
The three terms are exactly the lattice images of the three terms in Eq.~\eqref{eq:contH}: hopping
(kinetic) dressed by the link $U_n$ so it stays gauge invariant, the staggered mass, and the
electric energy $\propto L_n^2$.

\subsection{Dimensionless parameters --- decoding $ag=1$, $m/g=1$}
It is standard to rescale $H$ into a dimensionless form $W=\frac{2}{ag^2}H$ governed by two numbers
\begin{equation}
x \equiv \frac{1}{(ag)^2}\quad(\text{inverse coupling}),\qquad
\mu \equiv \frac{2m}{ag^2}\quad(\text{dimensionless mass}),
\label{eq:dimensionless}
\end{equation}
so that $W=x\sum_n[\chi^\dagger_nU_n\chi_{n+1}+\text{h.c.}] + \mu\sum_n(-1)^n\chi^\dagger_n\chi_n + \sum_n L_n^2$.
Your proposal's primary setting $ag=1,\ m/g=1$ therefore means $x=1$ and a mass comparable to the
coupling --- the intermediate-coupling regime where confinement physics is most visible.
(Conventions for the exact numerical prefactor of $\mu$ differ between references; the \emph{ratios}
$ag$ and $m/g$ are what is physically fixed.)

\begin{checkbox}{Checkpoint --- you understand stage 2 when you can answer:}
\begin{itemize}[leftmargin=1.4em,nosep]
\item Why does the hopping term carry a factor $U_n$ rather than being a bare $\chi^\dagger_n\chi_{n+1}$?
\item What physical quantity does $(-1)^n$ encode?
\item How do $x$ and $\mu$ relate to $ag$ and $m/g$?
\end{itemize}
\end{checkbox}

\section{Gauss's law and eliminating the gauge field}\label{sec:gauss}

\begin{whybox}
This is stage 3 and the conceptual crux of the project: it is \emph{why} $N=4$ sites need only
4 qubits instead of also needing qubits for the gauge field. \newx
\end{whybox}

\subsection{Gauss's law as a constraint}
On the lattice, Gauss's law (the generator of local gauge transformations) relates the jump in the
electric field across a site to the charge sitting there:
\begin{equation}
L_n - L_{n-1} = Q_n, \qquad Q_n = \chi^\dagger_n\chi_n - \frac{1-(-1)^n}{2}.
\label{eq:gausslaw}
\end{equation}
The subtraction is the staggered ``Dirac sea'': odd sites are filled in the vacuum, so the physical
charge is measured relative to that. Physical states must satisfy Eq.~\eqref{eq:gausslaw} for
\emph{every} $n$ --- it is a constraint, not a dynamical equation.

\subsection{Solving the constraint in 1+1D (the key trick)}
In one spatial dimension Eq.~\eqref{eq:gausslaw} is a simple recursion. With open boundary
conditions and a background field $L_0=\varepsilon$ fixed at the left edge, telescoping gives
\begin{equation}
\boxed{\; L_n = \varepsilon + \sum_{k=1}^{n} Q_k \;}
\label{eq:Lsolved}
\end{equation}
\emph{Every} link field is now an explicit function of the fermion charges. The gauge field is no
longer an independent degree of freedom --- this is the elimination. The constant $\varepsilon$ is
the \textbf{external field}: it is your proposal's post-quench parameter $L=5$.

\subsection{The price: a non-local interaction}
Substituting Eq.~\eqref{eq:Lsolved} into the electric energy $\sum_n L_n^2$ converts a local term
into a \emph{long-range} charge--charge interaction,
\begin{equation}
\sum_n L_n^2 = \sum_n\Bigl(\varepsilon + \sum_{k\le n} Q_k\Bigr)^2
= \underbrace{\sum_n\varepsilon^2}_{\text{const}}
+ 2\varepsilon\sum_n\!\sum_{k\le n}\!Q_k
+ \sum_{k\le k'} c_{kk'}\,Q_k Q_{k'} .
\label{eq:nonlocal}
\end{equation}
The double sum is a one-dimensional Coulomb potential: its strength grows with the separation
$|k-k'|$, which is precisely the \textbf{linear confining potential} of Sec.~\ref{sec:gauge}. So the
gauge field did not disappear --- it reappeared as a confining force between charges, and as a
\emph{linear coupling to $\varepsilon$} that is exactly how the external field $L=5$ enters the
dynamics.

\begin{remark}
Trade-off worth internalizing: eliminating the gauge field saves qubits (4 instead of $4+$links) but
makes $H$ \emph{non-local} (all-to-all charge interactions). In 1+1D this is a great deal; in higher
dimensions Gauss's law cannot be solved this cleanly, which is why people keep the gauge field
explicit there. Knowing this boundary is exactly the ``ingredient beyond course coverage'' your
proposal names.
\end{remark}

\section{Jordan--Wigner: fermions to a 4-qubit Pauli sum}\label{sec:jw}

\begin{whybox}
Stage 4. You have this from class (\known), so the goal here is just to connect it cleanly to the
output of stage 3 and to the conserved-charge symmetry your ansatz must respect.
\end{whybox}

The Jordan--Wigner map turns the 1D fermion chain into spins,
\begin{equation}
\chi_n = \Bigl(\prod_{l<n}(-Z_l)\Bigr)\sigma^-_n,\qquad
\chi^\dagger_n\chi_n = \frac{1+Z_n}{2},
\label{eq:jw}
\end{equation}
with $\sigma^\pm=(X\pm iY)/2$. Two consequences you will use directly:
\begin{itemize}[leftmargin=1.5em]
\item \textbf{Hopping becomes the tight-binding/XX form.} For nearest neighbors the JW strings
cancel and $\chi^\dagger_n\chi_{n+1}+\text{h.c.}\to \tfrac12(X_nX_{n+1}+Y_nY_{n+1})$ --- the same
tight-binding chain from class. (When the link $U_n$ has been eliminated as in Sec.~\ref{sec:gauss},
this is exactly what you get; otherwise $U_n$ rides along.)
\item \textbf{Mass and interaction terms become $Z$'s.} $\chi^\dagger_n\chi_n\to(1+Z_n)/2$, so the
staggered mass and the Coulomb term in Eq.~\eqref{eq:nonlocal} become sums of $Z_n$ and $Z_kZ_{k'}$.
\end{itemize}
The result is a 4-qubit Pauli sum $H=\sum_j c_j P_j$. The \emph{total charge} $\sum_n Q_n$
is conserved, which in spin language is the total magnetization $\sum_n Z_n$: your ansatz should stay
in one magnetization sector (stage 5).

\begin{checkbox}{Checkpoint --- stages 2--4 cohere when you can:}
write down, on paper for $N=4$, the chain $Q_1,\dots,Q_4$, use Eq.~\eqref{eq:Lsolved} to get
$L_1,\dots,L_4$ in terms of the $Q$'s and $\varepsilon$, and identify which Pauli strings appear
after JW. Reproducing this by hand is your Week-1 ``verify against ED'' deliverable.
\end{checkbox}

% ==================================================================
\part{Dynamics: ground state, quench, and two ways to evolve}
% ==================================================================

\section{Ground state by VQE (brief --- you have this)}\label{sec:vqe}
The variational principle $\ev{\psi(\theta)|H|\psi(\theta)}\ge E_0$ underwrites VQE: prepare a
parameterized state $\ket{\psi(\theta)}$ with a hardware-efficient ansatz (layers of $R_y,R_z$
rotations interleaved with linear-chain CNOTs, per your proposal), and minimize the energy with a
classical optimizer (Adam, $\|\nabla E\|<10^{-4}$). Respect the conserved magnetization sector so
you do not optimize into an unphysical charge state. Your success metric is the energy-accuracy ratio
\begin{equation}
r(E)=\frac{E_{\text{VQE}}-E_{\max}}{E_{\text{exact}}-E_{\max}}\in(-\infty,1],
\qquad r(E)=1 \text{ is exact},
\label{eq:rE}
\end{equation}
targeting $r(E)>0.99$ (paper Fig.~5). $E_{\max}$ is the largest eigenvalue, used to normalize so the
metric is scale-free. Nothing here is new; the only project-specific subtlety is keeping the ansatz
charge-conserving.

\section{The quench and Trotter baseline}\label{sec:trotter}

\subsection{What a quench is}
A \emph{quantum quench} means: prepare the system in the ground state of one Hamiltonian, then at
$t=0$ \emph{suddenly} change a parameter and let it evolve under the new Hamiltonian. Here you prepare
the ground state (stage 5), then switch on the external field $\varepsilon=L=5$ and evolve
$\ket{\psi(t)}=e^{-iHt}\ket{\psi_0}$. Because $\ket{\psi_0}$ is \emph{not} an eigenstate of the new
$H$, observables evolve nontrivially --- this is out-of-equilibrium dynamics, and it is what drives
Schwinger pair production.

\subsection{Second-order Suzuki--Trotter}
Split $H=\sum_j H_j$ into easily-exponentiated pieces (hopping, mass, Coulomb). The symmetric
(second-order) product formula for one step $\delta t$ is
\begin{equation}
S_2(\delta t)=\Bigl(\prod_{j} e^{-iH_j \delta t/2}\Bigr)\Bigl(\prod_{j\ \text{reversed}} e^{-iH_j \delta t/2}\Bigr),
\qquad e^{-iHt}\approx \big[S_2(\delta t)\big]^{r},\ \delta t=t/r,
\label{eq:suzuki}
\end{equation}
with local error $O(\delta t^3)$ per step and global error $O(\delta t^2)$. This is your
course-technique baseline against which the variational method is compared.

\begin{remark}[Where the folder's Bosse paper fits]
The paper actually sitting in your folder (time-dependent product formulas) is about making
Eq.~\eqref{eq:suzuki} \emph{cheaper} when $H$ has very different energy scales --- e.g.\ a ``large''
term (here the strong external field / Coulomb energy at $L=5$) and a ``small'' term (hopping). If
you do the stretch goal of varying $L$, that paper is the right reference for an improved Trotter
baseline. It is a bonus, not a core requirement.
\end{remark}

\section{McLachlan variational quantum simulation (the heart of the project)}\label{sec:mclachlan}

\begin{whybox}
Stage 7b. This is the paper's central method and the most important \newx{} topic. The idea: instead
of a deep Trotter circuit, keep a \emph{fixed shallow ansatz} $\ket{\psi(\theta(t))}$ and move the
parameters $\theta(t)$ so the state tracks true Schr\"odinger evolution as closely as the ansatz
allows.
\end{whybox}

\subsection{The time-dependent variational principle}
Exact evolution obeys $\partial_t\ket{\psi}=-iH\ket{\psi}$. Within the ansatz manifold, the state can
only move along directions $\ket{\partial_i\psi}\equiv\partial\ket{\psi}/\partial\theta_i$, so its
time derivative is forced to be $\ket{\dot\psi}=\sum_i \dot\theta_i\ket{\partial_i\psi}$. In general
this cannot equal $-iH\ket{\psi}$ exactly. \textbf{McLachlan's variational principle} chooses
$\dot\theta$ to minimize the residual
\begin{equation}
\delta\,\Bigl\| \bigl(\partial_t + iH\bigr)\ket{\psi(\theta)} \Bigr\| = 0 .
\label{eq:mclachlan}
\end{equation}

\subsection{Deriving the equations of motion}
Write the squared norm with $\ket{\dot\psi}=\sum_i\dot\theta_i\ket{\partial_i\psi}$:
\begin{align}
\Bigl\|(\partial_t+iH)\ket{\psi}\Bigr\|^2
&= \sum_{i,j}\dot\theta_i\dot\theta_j\braket{\partial_i\psi}{\partial_j\psi}
+ \ev{\psi|H^2|\psi} \notag\\
&\quad + i\sum_i\dot\theta_i\Bigl(\ev{\partial_i\psi|H|\psi}-\ev{\psi|H|\partial_i\psi}\Bigr).
\label{eq:normexpand}
\end{align}
Minimizing over each $\dot\theta_i$ (set $\partial/\partial\dot\theta_i=0$) gives a linear system
\begin{equation}
\boxed{\ \sum_j A_{ij}\,\dot\theta_j = C_i\ },\qquad
A_{ij}=\Real\braket{\partial_i\psi}{\partial_j\psi},\qquad
C_i=\Imag\ev{\partial_i\psi|H|\psi}.
\label{eq:AthetaC}
\end{equation}
These are exactly the $A$ and $C$ of your proposal. Solve for $\dot\theta$ at each timestep and
integrate (forward Euler or Runge--Kutta) to march $\theta(t)$ forward.

\begin{keyidea}[Geometric meaning of $A$]
$A_{ij}=\Real\braket{\partial_i\psi}{\partial_j\psi}$ is the real part of the \emph{quantum geometric
tensor} --- the Fubini--Study metric on the space of quantum states, restricted to the ansatz
manifold. It is the \emph{same} object as the metric in quantum natural gradient. Equation
\eqref{eq:AthetaC} is therefore a geometric statement: it \emph{projects} the exact velocity
$-iH\ket{\psi}$ onto the tangent space of the manifold the ansatz can reach, using $A$ as the metric.
That is why the method is only as good as the ansatz's expressiveness.
\end{keyidea}

\subsection{Measuring $A$ and $C$: Hadamard-test circuits}
For an ansatz built from Pauli rotations $\ket{\psi}=\prod_k e^{-i\theta_k P_k/2}\ket{0}$, the
derivative inserts a generator: $\ket{\partial_i\psi}$ contains a factor $-\tfrac{i}{2}P_i$. Hence
both $A_{ij}=\Real\braket{\partial_i\psi}{\partial_j\psi}$ and $C_i=\Imag\ev{\partial_i\psi|H|\psi}$
reduce to overlaps of the form $\Real/\Imag\,\ev{0|U^\dagger P_i \cdots P_j U|0}$. These are measured with
the \textbf{Hadamard test}: an ancilla is put in superposition, controls the insertion of the Pauli
generators $P_i$ (and the Pauli terms of $H$), and the ancilla's $\ev{X}$ gives the real part,
$\ev{Y}$ the imaginary part:
\begin{equation}
\ev{X}_{\text{anc}}=\Real\ev{\psi|V|\psi},\qquad \ev{Y}_{\text{anc}}=\Imag\ev{\psi|V|\psi}.
\label{eq:hadamard}
\end{equation}
You build one auxiliary circuit per matrix element. Debugging these circuits is exactly the Week-4
task in your timeline.

\begin{remark}[Why bother, if Trotter already works?]
On a 4-qubit simulator both methods work; the \emph{point} of reproducing the variational method is
that it keeps the circuit \emph{shallow and fixed-depth} regardless of evolution time, trading circuit
depth for extra measurements (the $A,C$ estimation). That is the near-term-hardware advantage the
paper is demonstrating. Your fidelity comparison (next section) is what quantifies the price you pay
for that shallowness.
\end{remark}

\section{Observables and validation}\label{sec:obs}
Three numbers close the loop:
\begin{itemize}[leftmargin=1.5em]
\item $\ev{L(t)}$ --- the electric field, computed from Eq.~\eqref{eq:Lsolved} as an expectation of a
$Z$-Pauli sum. Watch it respond to the quench.
\item $\ev{n(t)}$ --- particle density, $\propto\sum_n\ev{\chi^\dagger_n\chi_n}$ mapped via
Eq.~\eqref{eq:jw}; the signature of pair production.
\item Fidelity $F(t)=|\braket{\psi_{\text{exact}}(t)}{\psi_{\text{VQS}}(t)}|^2$ against the exact
state from \texttt{scipy.linalg.expm} on the $2^4=16$-dimensional space (exact at $N=4$). This is your
ground truth (paper Fig.~6).
\end{itemize}

% ==================================================================
\part{Your learning roadmap}
% ==================================================================

\section{An ordered study sequence (mapped to your 5 weeks)}
Study in \emph{dependency order}, not pipeline order --- foundations first.

\begin{center}
\renewcommand{\arraystretch}{1.3}
\begin{longtable}{c >{\raggedright\arraybackslash}p{0.30\textwidth} >{\raggedright\arraybackslash}p{0.44\textwidth}}
\toprule
\textbf{Order} & \textbf{Learn this} & \textbf{Goal / deliverable} \\
\midrule\endhead
1 & Gauge theory crash course (Sec.~\ref{sec:gauge}) & Explain in one paragraph why local $U(1)$ symmetry forces a gauge field. \\
2 & Schwinger model phenomenology (Sec.~\ref{sec:gauge}) & Predict qualitatively what $\ev{n(t)},\ev{L(t)}$ should do after the quench. \\
3 & Kogut--Susskind \& dimensionless params (Sec.~\ref{sec:ks}) & Write Eq.~\eqref{eq:ksH}; convert $ag,m/g$ to $x,\mu$. \\
4 & Gauss-law elimination (Sec.~\ref{sec:gauss}) & Derive Eq.~\eqref{eq:Lsolved}; see the non-local term appear. \emph{(Week 1)} \\
5 & Jordan--Wigner assembly (Sec.~\ref{sec:jw}) & Build the explicit $N{=}4$ Pauli sum; verify vs.\ ED. \emph{(Week 1)} \\
6 & VQE specifics (Sec.~\ref{sec:vqe}) & Reproduce $r(E)>0.99$. \emph{(Week 2)} \\
7 & Quench + Suzuki--Trotter (Sec.~\ref{sec:trotter}) & Baseline $\ev{L},\ev{n}$ vs.\ time. \emph{(Week 3)} \\
8 & McLachlan TDVP + QGT (Sec.~\ref{sec:mclachlan}) & Derive Eq.~\eqref{eq:AthetaC} from scratch. \emph{(Week 4)} \\
9 & Hadamard-test estimation (Sec.~\ref{sec:mclachlan}) & Build $A,C$ circuits; debug. \emph{(Week 4)} \\
10 & Observables \& fidelity (Sec.~\ref{sec:obs}) & Reproduce Fig.~6; write up. \emph{(Week 5)} \\
\bottomrule
\end{longtable}
\end{center}

\section{Self-check questions (use these as gates)}
Do not move to the next module until you can answer the current one without notes.
\begin{enumerate}[leftmargin=1.6em]
\item Why does demanding \emph{local} (not global) $U(1)$ invariance require introducing $A_\mu$?
\item In 1+1D, why can the electric field be solved for algebraically instead of being an independent
quantum variable?
\item Starting from Gauss's law Eq.~\eqref{eq:gausslaw}, derive $L_n=\varepsilon+\sum_{k\le n}Q_k$.
Where does the external field $L=5$ enter?
\item After eliminating the gauge field, $H$ becomes non-local. What physical phenomenon is this
non-locality (the Coulomb term) responsible for?
\item Derive $A_{ij}\dot\theta_j=C_i$ from McLachlan's principle. What is the geometric meaning of
$A$?
\item Why is the variational (McLachlan) method potentially advantageous over Trotter on near-term
hardware, even though both reproduce the $N{=}4$ dynamics?
\item How does a Hadamard test return the real vs.\ imaginary part of a state overlap?
\end{enumerate}

\section{Curated resources}
Prioritized; start at the top of each block.

\paragraph{The source paper (read first, in layers).}
Nagano, Bapat, Bauer, \emph{Phys.\ Rev.\ D} \textbf{108}, 034501 (2023), arXiv:2302.10933. First pass:
Sec.~2 (Hamiltonian) and the figures you must reproduce (Figs.~5--6). Second pass: Sec.~3 (ansatz and
variational evolution). Download this --- it is not the PDF currently in your folder.

\paragraph{Gauge theory / Schwinger model (Modules 1--2).}
\begin{itemize}[nosep,leftmargin=1.4em]
\item Kogut \& Susskind, \emph{Phys.\ Rev.\ D} \textbf{11}, 395 (1975) --- the original staggered
lattice formulation (your reference [2]). Read for the Hamiltonian structure, skim the rest.
\item Banks, Susskind, Kogut, \emph{Phys.\ Rev.\ D} \textbf{13}, 1043 (1976) --- strong-coupling
analysis (your reference [3]).
\item For a gentle modern entry: any ``lattice gauge theory for quantum simulation'' review (e.g.\
Ba\~nuls et al., \emph{Eur.\ Phys.\ J.\ D} 2020, ``Simulating lattice gauge theories within quantum
technologies'') --- read the Schwinger-model sections.
\end{itemize}

\paragraph{Variational real-time evolution (Modules 8--9).}
\begin{itemize}[nosep,leftmargin=1.4em]
\item Yuan, Endo, Zhao, Li, Benjamin, \emph{Quantum} \textbf{3}, 191 (2019), ``Theory of variational
quantum simulation'' (your reference [4]) --- the definitive derivation of McLachlan's principle and
the $A,C$ matrices, including the Hadamard-test circuits. This is your primary technical reference for
Sec.~\ref{sec:mclachlan}.
\item Stokes, Izaac, Killoran, Carleo, \emph{Quantum} \textbf{4}, 269 (2020), ``Quantum natural
gradient'' --- for the quantum geometric tensor / Fubini--Study metric viewpoint of $A$.
\end{itemize}

\paragraph{VQE and Trotter (Modules 6--7, refreshers).}
\begin{itemize}[nosep,leftmargin=1.4em]
\item Peruzzo et al., \emph{Nat.\ Commun.} \textbf{5}, 4213 (2014) --- the original VQE (your
reference [5]).
\item PennyLane demos: ``A brief overview of VQE,'' ``Variational quantum time evolution,'' and the
Hadamard-test demo --- these mirror your \texttt{default.qubit} implementation almost line for line.
\item Bosse, Childs et al., \emph{Nat.\ Commun.} 2025 (the PDF in your folder) --- only if you pursue
the varying-$L$ stretch goal.
\end{itemize}

\appendix
\newpage
\section{Key-equation cheat sheet}
\begin{center}
\renewcommand{\arraystretch}{1.5}
\begin{tabular}{>{\raggedright\arraybackslash}p{0.30\textwidth} l}
\toprule
\textbf{Quantity} & \textbf{Expression} \\
\midrule
Covariant derivative & $D_\mu=\partial_\mu+igA_\mu$ \\
Kogut--Susskind $H$ & Eq.~\eqref{eq:ksH} \\
Dimensionless params & $x=1/(ag)^2,\quad \mu=2m/(ag^2)$ \\
Gauss's law & $L_n-L_{n-1}=Q_n$ \\
Eliminated field & $L_n=\varepsilon+\sum_{k\le n}Q_k$ \\
Jordan--Wigner number & $\chi^\dagger_n\chi_n=(1+Z_n)/2$ \\
Hopping $\to$ spins & $\chi^\dagger_n\chi_{n+1}+\text{h.c.}\to\tfrac12(X_nX_{n+1}+Y_nY_{n+1})$ \\
VQE metric & $r(E)=(E_{\text{VQE}}-E_{\max})/(E_{\text{exact}}-E_{\max})$ \\
2nd-order Trotter & $S_2(\delta t)=\prod_j e^{-iH_j\delta t/2}\prod_{j\,\text{rev}}e^{-iH_j\delta t/2}$ \\
McLachlan EOM & $\sum_j A_{ij}\dot\theta_j=C_i$ \\
$A$ matrix & $A_{ij}=\Real\braket{\partial_i\psi}{\partial_j\psi}$ \\
$C$ vector & $C_i=\Imag\ev{\partial_i\psi|H|\psi}$ \\
Fidelity & $F(t)=|\braket{\psi_{\text{exact}}(t)}{\psi_{\text{VQS}}(t)}|^2$ \\
\bottomrule
\end{tabular}
\end{center}

\section{Notation conventions}
\begin{center}
\renewcommand{\arraystretch}{1.3}
\begin{tabular}{ll}
\toprule
\textbf{Symbol} & \textbf{Meaning} \\
\midrule
$\chi_n,\chi^\dagger_n$ & staggered fermion annihilation/creation at site $n$ \\
$U_n=e^{i\theta_n}$ & gauge link operator on link $n$ \\
$L_n$ & electric-field operator on link $n$ (integer eigenvalues) \\
$Q_n$ & staggered charge at site $n$ \\
$\varepsilon$ / $L$ & external (background) electric field; $=5$ post-quench \\
$a$ & lattice spacing; $g$ = gauge coupling; $m$ = fermion mass \\
$x,\mu$ & dimensionless coupling and mass, Eq.~\eqref{eq:dimensionless} \\
$X,Y,Z$ & single-qubit Pauli operators; $\sigma^\pm=(X\pm iY)/2$ \\
$\theta$ & variational parameters of the ansatz \\
$A_{ij},C_i$ & McLachlan matrix and vector, Eq.~\eqref{eq:AthetaC} \\
$F(t)$ & fidelity against exact diagonalization \\
\bottomrule
\end{tabular}
\end{center}

\section{Acronyms}
\begin{center}
\renewcommand{\arraystretch}{1.25}
\begin{tabular}{ll}
\toprule
QED & quantum electrodynamics \\
KS & Kogut--Susskind (staggered fermions) \\
JW & Jordan--Wigner \\
VQE & variational quantum eigensolver \\
VQD & variational quantum deflation \\
VQS & variational quantum simulation \\
TDVP & time-dependent variational principle \\
QGT & quantum geometric tensor \\
ED & exact diagonalization \\
\bottomrule
\end{tabular}
\end{center}

\end{document}
